\documentclass[11pt,a4paper]{article}

% ============================================================
% Paquetes
% ============================================================
\usepackage[utf8]{inputenc}
\usepackage[T1]{fontenc}
\usepackage[spanish,es-noindentfirst]{babel}
\usepackage{geometry}
\geometry{a4paper, margin=2.5cm}
\usepackage{xcolor}
\usepackage{listings}
\usepackage{longtable}
\usepackage{array}
\usepackage{booktabs}
\usepackage{fancyhdr}
\usepackage{hyperref}
\usepackage{enumitem}
\usepackage{caption}
\usepackage{amssymb}
\usepackage{ragged2e}
\usepackage{seqsplit}

% ============================================================
% Colores (estados del checklist)
% ============================================================
\definecolor{cumpleColor}{RGB}{0,128,0}
\definecolor{parcialColor}{RGB}{184,134,11}
\definecolor{noaplicaColor}{RGB}{110,110,110}
\definecolor{fueraColor}{RGB}{31,90,160}
\definecolor{nocumpleColor}{RGB}{178,34,34}
\definecolor{codebg}{RGB}{245,245,245}
\definecolor{coderule}{RGB}{180,180,180}
\definecolor{linknavy}{RGB}{20,60,120}
\definecolor{headergray}{RGB}{90,90,90}

\newcommand{\Cumple}{\textcolor{cumpleColor}{\textbf{Cumple}}}
\newcommand{\Parcial}{\textcolor{parcialColor}{\textbf{Parcial}}}
\newcommand{\NoAplica}{\textcolor{noaplicaColor}{No aplica}}
\newcommand{\Fuera}{\textcolor{fueraColor}{Fuera de alcance}}
\newcommand{\NoCumple}{\textcolor{nocumpleColor}{\textbf{No cumple}}}

% ============================================================
% hyperref
% ============================================================
\hypersetup{
    colorlinks=true,
    linkcolor=linknavy,
    citecolor=linknavy,
    urlcolor=linknavy,
    pdftitle={Validacion de Seguridad de Entradas de Datos - AprendIA},
    pdfauthor={Equipo AprendIA}
}

% ============================================================
% Encabezado / pie de pagina
% ============================================================
\pagestyle{fancy}
\fancyhf{}
\renewcommand{\headrulewidth}{0pt}
\lfoot{\footnotesize\textcolor{headergray}{Seguridad de la Información \textbullet\ Equipo AprendIA}}
\rfoot{\footnotesize\textcolor{headergray}{\thepage}}

% ============================================================
% listings: estilos de código
% ============================================================
\renewcommand{\lstlistingname}{Listado}
\captionsetup[lstlisting]{font=small,labelfont=bf}

\lstdefinestyle{javastyle}{
    language=Java,
    backgroundcolor=\color{codebg},
    basicstyle=\ttfamily\scriptsize,
    numbers=left,
    numberstyle=\tiny\color{gray},
    numbersep=8pt,
    frame=single,
    framerule=0.3pt,
    rulecolor=\color{coderule},
    breaklines=true,
    breakatwhitespace=false,
    keywordstyle=\color{blue!55!black}\bfseries,
    commentstyle=\color{green!35!black}\itshape,
    stringstyle=\color{red!55!black},
    showstringspaces=false,
    tabsize=2,
    xleftmargin=1em,
    xrightmargin=1em
}

\lstdefinestyle{dartstyle}{
    backgroundcolor=\color{codebg},
    basicstyle=\ttfamily\scriptsize,
    numbers=left,
    numberstyle=\tiny\color{gray},
    numbersep=8pt,
    frame=single,
    framerule=0.3pt,
    rulecolor=\color{coderule},
    breaklines=true,
    breakatwhitespace=false,
    showstringspaces=false,
    tabsize=2,
    xleftmargin=1em,
    xrightmargin=1em
}

\lstdefinestyle{plainstyle}{
    backgroundcolor=\color{codebg},
    basicstyle=\ttfamily\scriptsize,
    frame=single,
    framerule=0.3pt,
    rulecolor=\color{coderule},
    breaklines=true,
    breakatwhitespace=false,
    showstringspaces=false,
    tabsize=2,
    xleftmargin=1em,
    xrightmargin=1em
}

\lstdefinestyle{xmlstyle}{
    language=XML,
    backgroundcolor=\color{codebg},
    basicstyle=\ttfamily\scriptsize,
    numbers=left,
    numberstyle=\tiny\color{gray},
    numbersep=8pt,
    frame=single,
    framerule=0.3pt,
    rulecolor=\color{coderule},
    breaklines=true,
    breakatwhitespace=false,
    keywordstyle=\color{blue!55!black}\bfseries,
    commentstyle=\color{green!35!black}\itshape,
    stringstyle=\color{red!55!black},
    showstringspaces=false,
    xleftmargin=1em,
    xrightmargin=1em
}

\lstset{style=plainstyle}

% Placeholder para capturas de pantalla que residen en otro repositorio
\newcommand{\screenshotplaceholder}[3]{%
\begin{center}
\fbox{\parbox[c][#2][c]{#1}{\centering\footnotesize\itshape #3}}
\end{center}
}

\renewcommand{\contentsname}{Índice}

% ============================================================
\begin{document}
\shorthandoff{"}

% ============================================================
% Portada
% ============================================================
\begin{titlepage}
\centering
\vspace*{2cm}
{\Huge\bfseries Universidad Politecnica de Chiapas\par}
\vspace{0.4cm}
{\Large Ingenieria en Software\par}
\vspace{2cm}

\noindent\rule{\textwidth}{1pt}\par
\vspace{0.3cm}
{\Large\bfseries Validación de Seguridad de Entradas de Datos en la aplicación móvil AprendIA y su Backend\par}
\vspace{0.3cm}
\noindent\rule{\textwidth}{1pt}\par

\vspace{2.5cm}

\begin{tabular}{c}
\textbf{Materia:} Seguridad de la Información \\
\textbf{Profesor:} MC. Jose Alonso Macias Montoya \\
\textbf{Corte:} Corte 3 / Unidad III: Seguridad en Aplicaciones Móviles \\
\end{tabular}

\vspace{2cm}

{\bfseries\large Alumnos:\par}
\vspace{0.3cm}
Miguel Ángel Gutiérrez Gómez \\
Jesús Imanol Castillo Avendaño \\
Miguel Ángel Molina Gómez

\vspace{1cm}

{\bfseries\large Matrículas:\par}
\vspace{0.3cm}
233382 \\
233287 \\
233371

\vfill

{\bfseries Suchiapa, Chiapas, a 29 de julio de 2026}

\end{titlepage}

\newpage
\tableofcontents
\newpage

% ============================================================
\section{Introducción}

Este documento presenta la \textbf{validación de seguridad de las entradas de datos} de la aplicación educativa AprendIA, orientada a personas con discapacidad auditiva. El documento cubre dos componentes del sistema:

\begin{itemize}[leftmargin=1.5cm]
    \item El \textbf{cliente móvil} (Flutter/Dart, arquitectura limpia \texttt{\seqsplit{data/domain/presentation}}), evaluado originalmente contra el checklist de validación y sanitización de entradas.
    \item El \textbf{backend} (Spring Boot 3.4 / Java 21, \texttt{\seqsplit{aprendia-backend}}), que expone la API REST consumida por el cliente y por el panel administrativo, y que es responsable de toda la validación del lado del servidor que en la primera fase del proyecto había quedado documentada como \emph{fuera de alcance}.

\end{itemize}

El rasgo distintivo del cliente es que \textbf{no usa formularios de usuario/contraseña ni campos de texto libre} para el flujo del estudiante: la autenticación del alumno es por código QR con token, y los ejercicios se resuelven por selección de letras/imágenes, cámara y micrófono. El backend, en cambio, sí expone formularios estructurados (JSON) para autenticación de administradores, autorregistro de usuarios, registro de estudiantes (con datos personales, dirección y familiares) y administración de catálogos (dependencias, metodologías, municipios, perfiles, roles, estados).

\textbf{Trazabilidad.} El trabajo del cliente está implementado y versionado en la rama \texttt{\seqsplit{input-validate}} del repositorio móvil y es verificable con \texttt{\seqsplit{flutter analyze}}, \texttt{\seqsplit{flutter test}} y los archivos citados. El trabajo del backend está implementado en el repositorio \texttt{\seqsplit{aprendia-backend}} (rama \texttt{\seqsplit{somer}}) y es verificable revisando las anotaciones de \emph{Jakarta Bean Validation} en los DTO, el \texttt{\seqsplit{GlobalExceptionHandler}} y la configuración de \texttt{\seqsplit{SecurityConfig}}.

\section{Recursos Utilizados}

\begin{itemize}[leftmargin=1.5cm]
    \item \textbf{Cliente:} Flutter / Dart (SDK $\geq$3.4.4), arquitectura limpia; \texttt{\seqsplit{dio}} (red), \texttt{\seqsplit{flutter\_secure\_storage}} (sesión cifrada); análisis estático con \texttt{\seqsplit{flutter\_lints}}; pruebas con \texttt{\seqsplit{flutter\_test}}; CI con GitHub Actions.
    \item \textbf{Backend:} Spring Boot 3.4.0 / Java 21, Maven; \texttt{\seqsplit{spring-boot-starter-validation}} (Jakarta Bean Validation / Hibernate Validator); \texttt{\seqsplit{spring-boot-starter-security}} (JWT + roles); \texttt{\seqsplit{spring-boot-starter-data-jpa}} + PostgreSQL (consultas parametrizadas); \texttt{\seqsplit{springdoc-openapi}} (documentación generada desde las mismas anotaciones de validación); Flyway (migraciones versionadas).
    \item \textbf{Estándar de referencia:} OWASP \emph{MASVS} y \emph{Mobile Top 10} (cliente); OWASP \emph{ASVS} / \emph{API Security Top 10} (backend).
\end{itemize}

\section{Metodología}

Se aplicó la misma metodología en ambos componentes: (1) \textbf{mapeo de puntos de entrada} recorriendo el código fuente (\texttt{\seqsplit{lib/}} en el cliente, \texttt{\seqsplit{src/main/java}} en el backend); (2) \textbf{evaluación apartado por apartado} del checklist institucional, con estado (\Cumple/\Parcial/\NoAplica/\Fuera) y evidencia en código; (3) \textbf{verificación} —en el cliente con \texttt{\seqsplit{flutter analyze}} y \texttt{\seqsplit{flutter test}}; en el backend por inspección directa de las anotaciones de validación, del manejador global de excepciones y de la configuración de seguridad, dado que el proyecto aún no cuenta con una suite de pruebas automatizadas (ver Sección~\ref{sec:backend-parcial}).

\textbf{Criterio de estados}
\begin{itemize}[leftmargin=1.5cm]
    \item \Cumple\ --- implementado y verificado.
    \item \Parcial\ --- implementado pero pendiente de un paso o incompleto para algunos campos.
    \item \NoAplica\ --- no relevante para este componente (se justifica).
    \item \Fuera\ --- corresponde a otro componente del sistema (se justifica).
\end{itemize}

\section{Alcance: por qué no se valida lo externo al cliente}

\subsection{Autenticidad e integridad del token (backend) --- \Fuera}
\label{ssec:alcance-token}

Verificar firma, expiración, emisor y permisos exige material criptográfico y estado que solo tiene el servidor (principio \emph{never trust the client}). El cliente aporta defensa en profundidad: valida el \emph{formato} del QR, usa HTTPS en el login, manda el token en el cuerpo (no en la URL) y no lo registra en logs. Esta validación se documenta a detalle, con código, en la Sección~\ref{sec:backend-servidor}.

\subsection{Transporte cifrado (TLS/HTTPS) --- depende de infraestructura}

Los servicios de evaluación (\texttt{\seqsplit{20.245.212.209:5000/5001}}) exponen HTTP sin TLS: son IPs sin dominio ni certificado. Forzar \texttt{\seqsplit{https}} en el cliente antes de que el servidor ofrezca TLS dejaría la app sin conexión. \textbf{Riesgo residual: Alto} (MASVS-NETWORK); se escala como dependencia bloqueante a infraestructura. Este hallazgo sigue vigente y es independiente del trabajo de validación de entradas documentado en este informe.

\section{Tipos de entrada de datos y su validación (Cliente)}

La app tiene siete tipos de entrada de datos externos. La siguiente tabla mapea cada uno con su forma de validación, la \textbf{llamada a función} responsable y el archivo. Los bloques de código de cada validación se muestran en la Sección~\ref{sec:cliente-desarrollo}.

\begin{longtable}{p{3.1cm}p{6.6cm}p{4.3cm}}
\toprule
\textbf{Tipo de entrada} & \textbf{Forma de validación y llamada a función} & \textbf{Archivo} \\
\midrule
\endhead
Token de QR (texto) & Allowlist ASCII + longitud 8--4096: \texttt{\seqsplit{Validators.isValidQrToken(qrCode)}} & \texttt{\seqsplit{validators.dart}} \\
Respuesta de red (JSON) & Parseo defensivo con casts seguros: \texttt{\seqsplit{UserInfo.fromJson(json)}} & \texttt{\seqsplit{user\_info\_entity.dart}} \\
Foto de cámara (imagen) & Foto y referencia no vacías: \texttt{\seqsplit{ExerciseInputValidator.validateWritingInput(...)}} & \texttt{\seqsplit{evaluate\_writing\_exercise\_usecase.dart}} \\
Audio de micrófono & Audio y frase objetivo presentes: \texttt{\seqsplit{ExerciseInputValidator.validateReadingInput(...)}} & \texttt{\seqsplit{evaluate\_reading\_exercise\_usecase.dart}} \\
Selección (índice/opción) & Rango/pertenencia al conjunto: \texttt{\seqsplit{ExerciseInputValidator.validateSelectedOption(i, n)}} & \texttt{\seqsplit{exercise\_input\_validator.dart}} \\
URL de medios (remota) & Canonicalización con \texttt{\seqsplit{Uri.tryParse}}: \texttt{\seqsplit{sanitizeMediaUrl(url)}} & \texttt{\seqsplit{media\_url.dart}} \\
Sesión (token + datos) & Almacenamiento cifrado: \texttt{\seqsplit{SessionManager.saveSession(...)}} & \texttt{\seqsplit{session\_manager.dart}} \\
\bottomrule
\end{longtable}

\begin{lstlisting}[style=dartstyle,caption={Validación de la foto de cámara antes de evaluar (evaluate\_writing\_exercise\_usecase.dart)}]
Future<WritingFeedbackEntity> call({
  required File studentImage,
  required String originalImageUrl,
}) async {
  // Regla de negocio: no evaluar sin foto tomada ni imagen de referencia.
  if (!ExerciseInputValidator.validateWritingInput(
    photoPath: studentImage.path,
    referenceImage: originalImageUrl,
  )) {
    throw ArgumentError(
        'Entrada de escritura invalida: falta la foto o la imagen de referencia.');
  }
  return await _exerciseRepository.evaluateWritingExercise(
    studentImage: studentImage,
    originalImageUrl: originalImageUrl,
  );
}
\end{lstlisting}

\section{Librerías y frameworks: origen y uso (Cliente)}

Todas las dependencias provienen del repositorio oficial de paquetes de Dart/Flutter, \textbf{pub.dev}. La validación de entradas se implementa con un \textbf{módulo propio} (\texttt{\seqsplit{Validators}}), \emph{sin} una librería externa de validación: decisión de diseño para tener un único punto auditable y no depender de terceros en la lógica de seguridad.

\begin{longtable}{p{3.2cm}p{3.4cm}p{7.4cm}}
\toprule
\textbf{Librería (versión)} & \textbf{Origen} & \textbf{Uso en el manejo de entradas} \\
\midrule
\endhead
\texttt{\seqsplit{flutter\_secure\_storage}} (9.2.2) & pub.dev, publisher \texttt{\seqsplit{juliansteenbakker.nl}} & Guarda la sesión (token y datos del alumno, ya validados) en almacenamiento cifrado: Keystore (Android) / Keychain (iOS). \\
\texttt{\seqsplit{dio}} (5.8.0) & pub.dev, publisher \texttt{\seqsplit{dio.pub}} & Cliente HTTP. Serialización segura con \texttt{\seqsplit{FormData}}/\texttt{\seqsplit{Map}}/\texttt{\seqsplit{queryParameters}} (sin concatenar), \texttt{\seqsplit{ErrorInterceptor}} y \texttt{\seqsplit{LogInterceptor}} gateado sin cuerpos. \\
\texttt{\seqsplit{mobile\_scanner}} (5.0.1) & pub.dev, publisher \texttt{\seqsplit{juliansteenbakker.nl}} & Lee el valor crudo del QR (entrada externa) que después valida \texttt{\seqsplit{Validators.isValidQrToken}}. \\
\texttt{\seqsplit{equatable}} (2.0.5) & pub.dev, publisher \texttt{\seqsplit{felangel.dev}} & Igualdad por valor en las entidades de dominio (p. ej. \texttt{\seqsplit{UserInfo}}). \\
\texttt{\seqsplit{flutter\_lints}} (4.0.0, dev) & pub.dev, equipo Dart/Flutter & Reglas de análisis estático; gate de calidad en CI. \\
\bottomrule
\end{longtable}

\noindent\textbf{Dependencia retirada.} \texttt{\seqsplit{string\_validator}} (1.2.0) estaba declarada pero solo se usaba de forma cosmética (\texttt{\seqsplit{isNumeric}} para un tamaño de fuente). Se retiró y su uso se reemplazó por \texttt{\seqsplit{Validators.isNumeric}}, eliminando una dependencia innecesaria de la superficie de ataque.

\noindent\textbf{Herramientas de gobernanza (CI, no son dependencias de la app).} \texttt{\seqsplit{osv-scanner}} (Google) escanea dependencias vulnerables contra la base OSV; \texttt{\seqsplit{gitleaks}} detecta secretos; ambas corren en GitHub Actions (\texttt{\seqsplit{.github/workflows/security.yml}}).

\section{Desarrollo: apartados validados (Cliente, con código y evidencia)}
\label{sec:cliente-desarrollo}

\subsection{1. Validación del lado del cliente --- \Cumple}

\textbf{Estado:} \Cumple. La única entrada externa arbitraria (el token del QR) se valida por \textbf{allowlist} (ASCII imprimible, longitud 8--4096) \emph{antes} de enviarse. No valida autenticidad (eso es del backend): es defensa en profundidad.

\textbf{Dónde se cumple:}
\begin{itemize}[leftmargin=1.2cm,nosep]
\item \texttt{\seqsplit{lib/core/validation/validators.dart}}
\item \texttt{\seqsplit{lib/auth/presentation/viewmodels/qr\_scanner\_viewmodel.dart}}
\end{itemize}

\begin{lstlisting}[style=dartstyle,caption={Validación de formato del token de QR (validators.dart)}]
static final RegExp printableAscii = RegExp(r'^[\x21-\x7E]+$');

static bool isAllowedFreeText(String value,
    {int minLength = 1, int maxLength = 4096}) {
  if (!hasLengthBetween(value, minLength, maxLength)) return false;
  return printableAscii.hasMatch(value);
}

// Solo formato; la autenticidad la valida el backend.
static bool isValidQrToken(String value) =>
    isAllowedFreeText(value, minLength: 8, maxLength: 4096);
\end{lstlisting}

\begin{lstlisting}[style=dartstyle,caption={Rechazo del QR con formato inválido antes de llamar a la API (qr\_scanner\_viewmodel.dart)}]
if (!Validators.isValidQrToken(qrCode)) {
  _errorMessage = 'El codigo QR no tiene un formato valido.';
  _lastScannedCode = null; // permite reintentar
  notifyListeners();
  return;
}
\end{lstlisting}

\screenshotplaceholder{9cm}{5.5cm}{Cliente: banner de error ante un QR con formato inválido.\\(img/cap01\_qr\_invalido.png --- repositorio móvil)}
\captionof{figure}{Cliente: banner de error ante un QR con formato inválido.}

\subsection{3 y 8c. Validación de tipo / JSON (parseo defensivo) --- \Cumple}

\textbf{Estado:} \Cumple. El parseo de la respuesta de login no confía en la forma del JSON: accede por pasos con casts seguros y lanza una \texttt{\seqsplit{FormatException}} controlada si falta un campo obligatorio, en vez de un crash por desreferencia de \texttt{\seqsplit{null}}. Usa \texttt{\seqsplit{int.tryParse}} para el \texttt{\seqsplit{id}}.

\textbf{Dónde se cumple:}
\begin{itemize}[leftmargin=1.2cm,nosep]
\item \texttt{\seqsplit{user\_info\_entity.dart}}
\end{itemize}

\begin{lstlisting}[style=dartstyle,caption={fromJson defensivo (user\_info\_entity.dart)}]
factory UserInfo.fromJson(Map<String, dynamic> json) {
  final data = json['data'];
  if (data is! Map<String, dynamic>) {
    throw const FormatException('Respuesta de login sin "data" valido.');
  }
  final userInfo = data['userInfo'];
  if (userInfo is! Map<String, dynamic>) {
    throw const FormatException('Respuesta de login sin "userInfo" valido.');
  }
  final id = int.tryParse('${userInfo['studentId']}');
  final name = userInfo['nombre'];
  final token = data['token'];
  if (id == null || name is! String || token is! String) {
    throw const FormatException('Respuesta de login incompleta.');
  }
  // ... campos opcionales caen a cadena vacia ...
}
\end{lstlisting}

\subsection{4. Lógica de negocio --- \Cumple}

\textbf{Estado:} \Cumple. \texttt{\seqsplit{ExerciseInputValidator}} comprueba las invariantes de negocio \emph{antes} de evaluar: audio y frase para lectura; foto e imagen de referencia para escritura; índice de opción dentro del conjunto; tema no vacío. Se aplica en los use cases de evaluación.

\textbf{Dónde se cumple:}
\begin{itemize}[leftmargin=1.2cm,nosep]
\item \texttt{\seqsplit{exercises/domain/validation/exercise\_input\_validator.dart}}
\item \texttt{\seqsplit{evaluate\_reading\_exercise\_usecase.dart}}
\end{itemize}

\begin{lstlisting}[style=dartstyle,caption={Invariantes de negocio (exercise\_input\_validator.dart)}]
static bool validateReadingInput({
    required String audioPath, required String targetPhrase}) {
  if (audioPath.trim().isEmpty) return false;
  final phrase = targetPhrase.trim();
  return phrase.isNotEmpty && phrase.length <= maxTargetPhraseLength;
}

static bool validateSelectedOption(int index, int optionCount) =>
    Validators.isInRange(index, min: 0, max: optionCount - 1);
\end{lstlisting}

\subsection{7 y 8d. Validación contextual / limpieza --- \Cumple}

\textbf{Estado:} \Cumple. Normalización unificada para comparar respuestas: minúsculas, sin tildes ni diéresis, espacios colapsados, cubriendo formas Unicode precompuestas y descompuestas. Preserva la regla de negocio: la ñ \textbf{no} se degrada a n.

\textbf{Dónde se cumple:}
\begin{itemize}[leftmargin=1.2cm,nosep]
\item \texttt{\seqsplit{core/text/contextual\_normalizer.dart}}
\end{itemize}

\begin{lstlisting}[style=dartstyle,caption={Normalización de texto (contextual\_normalizer.dart)}]
static String normalize(String input) {
  var s = input.toLowerCase().trim();
  s = s.replaceAll(RegExp(r'\s+'), ' ');
  // Protege la nn antes de tocar diacriticos y la restaura al final.
  const enyeMarker = '#ENYE#';
  s = s.replaceAll('n~', enyeMarker);
  _accentMap.forEach((a, base) => s = s.replaceAll(a, base)); // a->a, u->u...
  s = s.replaceAll(_combiningMarks, ''); // formas descompuestas (NFD)
  return s.replaceAll(enyeMarker, 'n~');
}
\end{lstlisting}

\subsection{8b. Filtrado (allowlist) --- \Cumple}

\textbf{Estado:} \Cumple. Se usa \textbf{allowlist} (cerrada por defecto) en lugar de blacklist: una entrada libre es válida solo si encaja en \texttt{\seqsplit{printableAscii}} y en un rango de longitud. Reutilizable desde el módulo central.

\textbf{Dónde se cumple:}
\begin{itemize}[leftmargin=1.2cm,nosep]
\item \texttt{\seqsplit{validators.dart}}
\end{itemize}

\subsection{8f. Funciones/librerías seguras (serialización HTTP; ORM N/A) --- \Cumple}

\textbf{Estado:} \Cumple. Los paths interpolados usan \texttt{\seqsplit{Uri.encodeComponent}} y los cuerpos usan \texttt{\seqsplit{FormData}}/\texttt{\seqsplit{Map}}; no hay concatenación manual de URLs. Se blinda con una \textbf{prueba de arquitectura} que falla si reaparecen los patrones frágiles. \textbf{ORM = No aplica}: no hay base de datos local con consultas de usuario.

\textbf{Dónde se cumple:}
\begin{itemize}[leftmargin=1.2cm,nosep]
\item \texttt{\seqsplit{test/architecture/safe\_http\_serialization\_test.dart}}
\end{itemize}

\begin{lstlisting}[style=dartstyle,caption={Prueba de arquitectura que impide regresiones (safe\_http\_serialization\_test.dart)}]
test('no queda ningun split(":") de URL en lib/', () {
  final offenders = <String>[];
  for (final file in dartFiles) {
    if (file.readAsStringSync().contains("split(':')")) {
      offenders.add(file.path);
    }
  }
  expect(offenders, isEmpty);
});
\end{lstlisting}

\subsection{8h. Canonicalización / path traversal --- \Cumple}

\textbf{Estado:} \Cumple. \texttt{\seqsplit{sanitizeMediaUrl}} reemplaza un saneo frágil (\texttt{\seqsplit{split(':')}}) que perdía el esquema de la URL (\texttt{\seqsplit{https://x}} $\to$ \texttt{\seqsplit{//x}}). Usa \texttt{\seqsplit{Uri.tryParse}} como reja de validez.

\textbf{Dónde se cumple:}
\begin{itemize}[leftmargin=1.2cm,nosep]
\item \texttt{\seqsplit{core/net/media\_url.dart}}
\end{itemize}

\begin{lstlisting}[style=dartstyle,caption={Canonicalización de URLs de medios (media\_url.dart)}]
String sanitizeMediaUrl(String raw) {
  final trimmed = raw.trim();
  if (trimmed.isEmpty) return '';
  // Prefijo espurio antes de ':http' -> conservar desde 'http'.
  final marker = trimmed.indexOf(':http');
  final candidate = marker >= 0 ? trimmed.substring(marker + 1) : trimmed;
  if (Uri.tryParse(candidate) == null) return ''; // reja de validez
  return candidate;
}
\end{lstlisting}

\subsection{9. Uso de librerías/frameworks de validación --- \Cumple}

\textbf{Estado:} \Cumple. Módulo \texttt{\seqsplit{Validators}} como punto único de validación. Se \textbf{retiró} la dependencia \texttt{\seqsplit{string\_validator}}, que estaba declarada y solo se usaba de forma cosmética, reemplazándola por \texttt{\seqsplit{Validators.isNumeric}}.

\textbf{Dónde se cumple:}
\begin{itemize}[leftmargin=1.2cm,nosep]
\item \texttt{\seqsplit{validators.dart}}
\item \texttt{\seqsplit{pubspec.yaml}}
\end{itemize}

\subsection{11. Gestión de errores y almacenamiento seguro --- \Cumple}

\textbf{Estado:} \Cumple. Mensajes de interfaz genéricos y logs sin secretos: se eliminaron/gatearon todos los \texttt{\seqsplit{print()}}; ninguno emite token, cuerpos, cabeceras ni datos personales, y el \texttt{\seqsplit{LogInterceptor}} de Dio solo se registra en \texttt{\seqsplit{kDebugMode}} y sin cuerpos. Además, la sesión se guarda \textbf{cifrada} con \texttt{\seqsplit{flutter\_secure\_storage}} y \texttt{\seqsplit{clearSession()}} borra las 6 claves.

\textbf{Dónde se cumple:}
\begin{itemize}[leftmargin=1.2cm,nosep]
\item \texttt{\seqsplit{dio\_client.dart}}
\item \texttt{\seqsplit{session\_manager.dart}}
\end{itemize}

\begin{lstlisting}[style=dartstyle,caption={LogInterceptor gateado y sin cuerpos (dio\_client.dart)}]
if (kDebugMode) {
  _dio.interceptors.add(LogInterceptor(
    request: true, requestHeader: false,
    responseHeader: false, responseBody: false,
  ));
}
\end{lstlisting}

\begin{lstlisting}[style=dartstyle,caption={Sesión en almacenamiento cifrado (session\_manager.dart)}]
final FlutterSecureStorage _storage; // Keystore/Keychain

Future<void> clearSession() async {
  for (final key in _allKeys) { // las 6 claves
    await _storage.delete(key: key);
  }
}
\end{lstlisting}

\section{Apartados No aplica (Cliente, justificados)}

\begin{description}[leftmargin=1.2cm,style=nextline]
\item[5. Patrones (email/Luhn/contraseñas)] \textbf{Estado:} \NoAplica. \textbf{Por qué no aplica:} la app no captura correo, contraseñas ni tarjetas: autenticación por QR.
\item[6. Validación cruzada] \textbf{Estado:} \NoAplica. \textbf{Por qué no aplica:} no hay formularios multi-campo; el caso análogo (índice dentro de opciones) lo cubre \texttt{\seqsplit{validateSelectedOption}}.
\item[8a. Escapado HTML/JS/SQL] \textbf{Estado:} \NoAplica. \textbf{Por qué no aplica:} interfaz nativa; sin SQL de usuario expuesto.
\item[8e. Codificación] \textbf{Estado:} \NoAplica. \textbf{Por qué no aplica:} el token va en el cuerpo; se usa \texttt{\seqsplit{Uri.encodeComponent}} en segmentos de path.
\item[8g. Reemplazo de caracteres] \textbf{Estado:} \NoAplica. \textbf{Por qué no aplica:} no hay entrada de usuario con comillas/scripts.
\item[8i. Escape de salida] \textbf{Estado:} \NoAplica. \textbf{Por qué no aplica:} interfaz 100\,\% nativa Flutter.
\item[10. Educación/capacitación] \textbf{Estado:} \NoAplica. \textbf{Por qué no aplica:} no evidenciable desde el repo; \texttt{\seqsplit{CONTRIBUTING.md}} documenta las reglas y sirve de material de incorporación.
\end{description}

\section{Apartado Parcial (Cliente)}

\subsection*{8j. Revisiones y auditorías (SAST/pentest) --- \Parcial}

\textbf{Estado:} \Parcial. Se añadió un workflow de CI con \texttt{\seqsplit{flutter analyze}} como \emph{gate}, \texttt{\seqsplit{flutter test}}, \texttt{\seqsplit{osv-scanner}} (dependencias) y \texttt{\seqsplit{gitleaks}} (secretos), más una guía de contribución.

\textbf{Dónde se cumple:}
\begin{itemize}[leftmargin=1.2cm,nosep]
\item \texttt{\seqsplit{.github/workflows/security.yml}}
\item \texttt{\seqsplit{CONTRIBUTING.md}}
\end{itemize}

\textbf{Por qué es parcial:} el workflow debe ejecutarse en una PR real para confirmarse en producción, y el pentest manual está pendiente de agendar.

\section{Correcciones aplicadas (rama input-validate, Cliente)}

\begin{longtable}{p{4.1cm}p{2cm}p{6.9cm}}
\toprule
\textbf{Brecha} & \textbf{Severidad} & \textbf{Acción} \\
\midrule
\endhead
Token en \texttt{\seqsplit{SharedPreferences}} en claro; \texttt{\seqsplit{clearSession}} incompleto & Alta & \texttt{\seqsplit{flutter\_secure\_storage}}; \texttt{\seqsplit{clearSession}} borra las 6 claves \\
Token/datos en logs (\texttt{\seqsplit{print}} + \texttt{\seqsplit{LogInterceptor}} con cuerpos) & Alta & \texttt{\seqsplit{print}} eliminados/gateados; \texttt{\seqsplit{LogInterceptor}} solo en \texttt{\seqsplit{kDebugMode}} y sin cuerpos \\
QR sin validación de formato & Media & Allowlist por regex antes de enviarlo \\
\texttt{\seqsplit{fromJson}} sin null-safety & Media & Parseo defensivo con \texttt{\seqsplit{FormatException}} \\
Reglas de negocio sin validar & Media & \texttt{\seqsplit{ExerciseInputValidator}} en los use cases \\
Saneo de URL frágil (pierde esquema) & Baja & \texttt{\seqsplit{sanitizeMediaUrl}}; 9 copias reemplazadas \\
Normalización divergente (diéresis) & Baja & \texttt{\seqsplit{ContextualNormalizer}} compartido \\
Dependencia cosmética \texttt{\seqsplit{string\_validator}} & Baja & Retirada; módulo \texttt{\seqsplit{Validators}} \\
\bottomrule
\end{longtable}

\noindent\textbf{Verificación reproducible} (rama \texttt{\seqsplit{input-validate}}): \texttt{\seqsplit{flutter analyze}} $\rightarrow$ \textbf{0 errores}; \texttt{\seqsplit{flutter test}} $\rightarrow$ \textbf{22/22} pruebas de seguridad en verde (unitarias + arquitectura), y 169 pruebas del proyecto en verde; \texttt{\seqsplit{grep}} de \texttt{\seqsplit{print(}} y \texttt{\seqsplit{split(':')}} en \texttt{\seqsplit{lib/}} $\rightarrow$ \textbf{0} en código.

\section{Tabla Resumen de Cumplimiento (Cliente)}
\label{sec:tabla-cliente}

\begin{longtable}{p{0.5cm}p{4cm}p{2.6cm}p{6.2cm}}
\toprule
\textbf{\#} & \textbf{Apartado} & \textbf{Estado} & \textbf{Evidencia} \\
\midrule
\endhead
1 & Validación cliente & \Cumple & \texttt{\seqsplit{core/validation/validators.dart}} \\
2 & Validación servidor & \Fuera & Backend; HTTPS login en cliente \\
3 & Validación de tipo & \Cumple & \texttt{\seqsplit{user\_info\_entity.dart}} \\
4 & Lógica de negocio & \Cumple & \texttt{\seqsplit{exercise\_input\_validator.dart}} \\
5 & Patrones (email/Luhn) & \NoAplica & Login QR; sin campos sensibles \\
6 & Validación cruzada & \NoAplica & Sin formularios multi-campo \\
7 & Validación contextual & \Cumple & \texttt{\seqsplit{core/text/contextual\_normalizer.dart}} \\
8a & Escapado HTML/JS/SQL & \NoAplica & Sin WebView/SQL de usuario \\
8b & Filtrado (allowlist) & \Cumple & \texttt{\seqsplit{validators.dart}} \\
8c & Validación tipo JSON & \Cumple & \texttt{\seqsplit{user\_info\_entity.dart}} \\
8d & Limpieza (trim/normalize) & \Cumple & \texttt{\seqsplit{contextual\_normalizer.dart}} \\
8e & Codificación & \NoAplica & Token en cuerpo; \texttt{\seqsplit{Uri.encodeComponent}} \\
8f & Librerías seguras / ORM N/A & \Cumple & \texttt{\seqsplit{test/architecture/}} \\
8g & Reemplazo de caracteres & \NoAplica & Sin entrada con scripts \\
8h & Canonicalización & \Cumple & \texttt{\seqsplit{core/net/media\_url.dart}} \\
8i & Escape de salida & \NoAplica & Interfaz nativa \\
8j & Revisiones/auditorías & \Parcial & \texttt{\seqsplit{.github/workflows/security.yml}} \\
9 & Librerías de validación & \Cumple & \texttt{\seqsplit{validators.dart}}; \texttt{\seqsplit{pubspec.yaml}} \\
10 & Educación/capacitación & \NoAplica & \texttt{\seqsplit{CONTRIBUTING.md}} \\
11 & Gestión de errores & \Cumple & \texttt{\seqsplit{dio\_client.dart}}; \texttt{\seqsplit{session\_manager.dart}} \\
\bottomrule
\end{longtable}

\noindent\textbf{Conteo final (20 apartados, Cliente):} \Cumple\ = \textbf{11} (1, 3, 4, 7, 8b, 8c, 8d, 8f, 8h, 9, 11) \quad \Parcial\ = \textbf{1} (8j) \quad \NoAplica\ = \textbf{7} (5, 6, 8a, 8e, 8g, 8i, 10) \quad \Fuera\ = \textbf{1} (2) \quad \NoCumple\ = 0.

\section{Índice de capturas requeridas (Cliente)}

Cada captura referida en esta sección reside en \texttt{\seqsplit{docs/img/}} del repositorio móvil (fuera del alcance de este repositorio de backend) y se documenta aquí únicamente por trazabilidad con el informe original.

\begin{longtable}{p{1.8cm}p{7cm}p{4.7cm}}
\toprule
\textbf{Apartado} & \textbf{Qué debe mostrar} & \textbf{Archivo} \\
\midrule
\endhead
1 & Banner "El código QR no tiene un formato válido" & \texttt{\seqsplit{img/cap01\_qr\_invalido.png}} \\
3 & Login exitoso tras un QR válido & \texttt{\seqsplit{img/cap03\_login\_ok.png}} \\
4 & Botón "Enviar" deshabilitado sin la entrada requerida & \texttt{\seqsplit{img/cap04\_regla\_negocio.png}} \\
7 & Banco de letras sin tildes ("Deletrear mi nombre") & \texttt{\seqsplit{img/cap07\_deletreo.png}} \\
8f & Consola con las pruebas en verde & \texttt{\seqsplit{img/cap08\_tests\_verdes.png}} \\
8h & Ejercicios con imagen renderizando bien & \texttt{\seqsplit{img/cap09\_imagenes\_ok.png}} \\
11 & \texttt{\seqsplit{logcat}} en login sin token ni nombre & \texttt{\seqsplit{img/cap11\_log\_sin\_token.png}} \\
\bottomrule
\end{longtable}

% ============================================================
% ==================  PARTE BACKEND  ========================
% ============================================================
\newpage
\section{Backend: Validación de Seguridad de Entradas de Datos}
\label{sec:backend-servidor}

Esta sección documenta, apartado por apartado y con código real del repositorio \texttt{\seqsplit{aprendia-backend}} (rama \texttt{\seqsplit{somer}}), todo lo relacionado con la validación de entradas que se ejecuta \textbf{del lado del servidor} --- precisamente lo que la Sección~\ref{ssec:alcance-token} había dejado registrado como \Fuera\ por no ser responsabilidad del cliente.

\subsection{Alcance y arquitectura}

El backend es una API REST construida con \textbf{Spring Boot 3.4.0} sobre \textbf{Java 21}, organizada por \emph{features} (\texttt{\seqsplit{feature/auth}}, \texttt{\seqsplit{feature/user}}, \texttt{\seqsplit{feature/catalogs}}, \texttt{\seqsplit{feature/healthCheck}}) más un paquete transversal \texttt{\seqsplit{security}} (filtros de autenticación) y \texttt{\seqsplit{exception}} (manejo centralizado de errores). Todas las entradas externas llegan como \textbf{JSON sobre HTTP} (cuerpos de petición, variables de ruta, parámetros de consulta y cabeceras); no hay renderizado de HTML del lado del servidor ni acceso a sistema de archivos a partir de datos de usuario.

\subsection{Recursos utilizados (Backend)}

\begin{itemize}[leftmargin=1.5cm]
    \item \textbf{Validación declarativa:} \texttt{\seqsplit{spring-boot-starter-validation}} (Jakarta Bean Validation, implementada por Hibernate Validator), aplicada sobre \emph{records} de Java que representan los DTO de entrada.
    \item \textbf{Autenticación/autorización:} \texttt{\seqsplit{spring-boot-starter-security}} + \texttt{\seqsplit{jjwt}} (JSON Web Tokens firmados con HMAC-SHA256) + \texttt{\seqsplit{@PreAuthorize}} basado en roles.
    \item \textbf{Persistencia:} \texttt{\seqsplit{spring-boot-starter-data-jpa}} sobre PostgreSQL, con consultas 100\,\% parametrizadas (sin \texttt{\seqsplit{@Query}} nativo en todo el proyecto).
    \item \textbf{Manejo de errores:} \texttt{\seqsplit{GlobalExceptionHandler}} (\texttt{\seqsplit{@RestControllerAdvice}}) centraliza la traducción de excepciones de validación a respuestas HTTP consistentes.
    \item \textbf{Documentación generada:} \texttt{\seqsplit{springdoc-openapi}} construye el contrato OpenAPI/Swagger a partir de las mismas anotaciones de validación, evitando divergencia entre documentación y código.
\end{itemize}

\subsection{Metodología (Backend)}

Se recorrieron los cinco puntos de entrada de la aplicación (controladores REST, cuerpos JSON deserializados, variables de ruta, parámetros de consulta y cabeceras de autenticación), se identificó el mecanismo de validación asociado a cada uno y se contrastó contra el mismo checklist institucional usado para el cliente. A diferencia del cliente, el backend \textbf{no cuenta todavía} con una suite de pruebas automatizadas (\texttt{\seqsplit{src/test/java}} está vacío) ni con un pipeline de CI; esto se documenta explícitamente como una brecha en la Sección~\ref{sec:backend-parcial}.

\subsection{Mapa de puntos de entrada (Backend)}

\begin{longtable}{p{4.4cm}p{6cm}p{3cm}}
\toprule
\textbf{Punto de entrada} & \textbf{Evidencia} & \textbf{Tipo de dato} \\
\midrule
\endhead
Login admin (cuerpo JSON) & \texttt{\seqsplit{AuthController.loginAdmin}} / \texttt{\seqsplit{UserLoginDTO}} & Credenciales (texto) \\
Login estudiante (cuerpo JSON) & \texttt{\seqsplit{AuthController.loginStudent}} / \texttt{\seqsplit{StudentQrLoginDTO}} & Token QR (texto) \\
Validación de token (cuerpo JSON) & \texttt{\seqsplit{AuthController.validateToken}} / \texttt{\seqsplit{TokenValidationRequestDTO}} & Token JWT (texto) \\
Autorregistro de usuario (cuerpo JSON) & \texttt{\seqsplit{UserController.registerUser}} / \texttt{\seqsplit{RegisterRequest}} & Usuario / correo / contraseña \\
Registro de estudiante (JSON anidado) & \texttt{\seqsplit{UserController.registerStudent}} / \texttt{\seqsplit{StudentRequest}} & Persona + dirección + familiares \\
Catálogos administrativos (cuerpo JSON) & \texttt{\seqsplit{Dependency/Metodology/Municipality/Profile/Role/State RequestDTO}} & Texto / numérico / booleano \\
Identificadores de recurso (ruta) & \texttt{\seqsplit{@PathVariable Long id}}, \texttt{\seqsplit{@PathVariable Integer id}} & Numérico \\
Paginación (parámetros de consulta) & \texttt{\seqsplit{Pageable}}, \texttt{\seqsplit{@PageableDefault}} & Numérico \\
Token de sesión (cabecera) & \texttt{\seqsplit{Authorization: Bearer <jwt>}} / \texttt{\seqsplit{JwtAuthenticationFilter}} & Token (texto) \\
Clave de servicio (cabecera) & \texttt{\seqsplit{X-API-KEY}} / \texttt{\seqsplit{ApiKeyFilter}} & Clave (texto) \\
\bottomrule
\end{longtable}

\subsection{Tipos de entrada de datos y su validación (Backend)}
\label{ssec:backend-tipos}

\begin{longtable}{p{3.2cm}p{6.4cm}p{4.2cm}}
\toprule
\textbf{Tipo de entrada} & \textbf{Forma de validación y llamada a función} & \textbf{Archivo} \\
\midrule
\endhead
Credenciales de login admin & \texttt{\seqsplit{@NotBlank}} en \texttt{\seqsplit{UserLoginDTO}}, verificadas por \texttt{\seqsplit{AuthenticationManager.authenticate(...)}} & \texttt{\seqsplit{UserLoginDTO.java}}, \texttt{\seqsplit{AuthServiceImpl.java}} \\
Código QR de estudiante & \texttt{\seqsplit{@NotBlank}} en \texttt{\seqsplit{StudentQrLoginDTO}}; búsqueda parametrizada \texttt{\seqsplit{UserRepository.findByQrCode(qr)}} & \texttt{\seqsplit{StudentQrLoginDTO.java}} \\
Token JWT recibido & Verificación de firma HMAC-SHA256 y expiración: \texttt{\seqsplit{JwtUtil.validateToken(token, userDetails)}} & \texttt{\seqsplit{JwtUtil.java}}, \texttt{\seqsplit{JwtAuthenticationFilter.java}} \\
Registro de usuario (usuario / correo / contraseña) & \texttt{\seqsplit{@NotBlank}}, \texttt{\seqsplit{@Email}}, \texttt{\seqsplit{@Size(min,max)}} en \texttt{\seqsplit{RegisterRequest}} + \texttt{\seqsplit{@Valid}} en el controlador & \texttt{\seqsplit{RegisterRequest.java}}, \texttt{\seqsplit{UserController.java}} \\
Registro de estudiante (datos anidados) & Validación en cascada \texttt{\seqsplit{@Valid}} sobre \texttt{\seqsplit{PersonRequest}}, \texttt{\seqsplit{AddressRequest}} y \texttt{\seqsplit{List<RelativeRequest>}} & \texttt{\seqsplit{StudentRequest.java}} \\
Referencias a catálogos (perfil, municipio, estado) & Existencia verificada en BD: \texttt{\seqsplit{ensureExists(repository, id, nombre)}} & \texttt{\seqsplit{UserServiceImpl.java}} \\
Duplicados de negocio (usuario, correo, CURP) & \texttt{\seqsplit{UserRepository.existsByUsername/existsByEmail}}, \texttt{\seqsplit{PersonRepository.existsByCurp}} & \texttt{\seqsplit{UserServiceImpl.java}} \\
Identificadores de ruta (\texttt{\seqsplit{id}}) & Coerción de tipo de Spring MVC + \texttt{\seqsplit{JpaRepository.existsById(id)}} & Controladores REST \\
Cuerpo JSON malformado / tipo incorrecto & \texttt{\seqsplit{HttpMessageNotReadableException}} $\to$ \texttt{\seqsplit{GlobalExceptionHandler.handleHttpMessageNotReadable}} & \texttt{\seqsplit{GlobalExceptionHandler.java}} \\
Clave de servicio (cabecera \texttt{\seqsplit{X-API-KEY}}) & Presencia y existencia no-obsoleta: \texttt{\seqsplit{ApiKeyRepository.findByKeyHash(key)}} & \texttt{\seqsplit{ApiKeyFilter.java}} \\
\bottomrule
\end{longtable}

\begin{lstlisting}[style=javastyle,caption={DTO de registro con validación declarativa (RegisterRequest.java)}]
public record RegisterRequest(

    @NotBlank(message = "El nombre de usuario es obligatorio.")
    @Size(min = 3, max = 50, message = "El nombre de usuario debe tener entre 3 y 50 caracteres.")
    String username,

    @NotBlank(message = "El correo electronico es obligatorio.")
    @Email(message = "El formato de correo electronico no es valido.")
    @Size(max = 100, message = "El correo electronico no puede superar los 100 caracteres.")
    String email,

    @NotBlank(message = "La contrasena es obligatoria.")
    @Size(min = 6, max = 100, message = "La contrasena debe tener entre 6 y 100 caracteres.")
    String password
) {}
\end{lstlisting}

\begin{lstlisting}[style=javastyle,caption={Activación de la validación en el controlador (UserController.java)}]
@PostMapping("/register")
@SecurityRequirements
public ResponseEntity<UserResponse> registerUser(
        @Valid @RequestBody RegisterRequest registerRequest) {
    UserResponse response = userService.registerUser(registerRequest);
    return new ResponseEntity<>(response, HttpStatus.CREATED);
}
\end{lstlisting}

\subsection{Librerías y frameworks: origen y uso (Backend)}

Todas las dependencias provienen del repositorio oficial \textbf{Maven Central}, resueltas por Maven a partir de \texttt{\seqsplit{pom.xml}} (padre \texttt{\seqsplit{spring-boot-starter-parent:3.4.0}}). No existe una dependencia declarada explícitamente para \texttt{\seqsplit{hibernate-validator}} ni \texttt{\seqsplit{jakarta.validation-api}}: ambas llegan de forma \textbf{transitiva} a través de \texttt{\seqsplit{spring-boot-starter-validation}}.

\begin{longtable}{p{4cm}p{3.2cm}p{6.8cm}}
\toprule
\textbf{Librería (versión)} & \textbf{Origen} & \textbf{Uso en el manejo de entradas} \\
\midrule
\endhead
\texttt{\seqsplit{spring-boot-starter-\allowbreak validation}} (3.4.0) & Maven Central, \texttt{\seqsplit{org.springframework.boot}} & Habilita \texttt{\seqsplit{@Valid}} y trae transitivamente Hibernate Validator (implementación de referencia de Jakarta Bean Validation) que ejecuta \texttt{\seqsplit{@NotNull}}/\texttt{\seqsplit{@NotBlank}}/\texttt{\seqsplit{@Size}}/\texttt{\seqsplit{@Min}}/\texttt{\seqsplit{@Email}} sobre los DTO. \\
\texttt{\seqsplit{spring-boot-starter-\allowbreak security}} (3.4.0) & Maven Central, \texttt{\seqsplit{org.springframework.security}} & Autenticación/autorización; \texttt{\seqsplit{BCryptPasswordEncoder}} para contraseñas; \texttt{\seqsplit{@PreAuthorize("hasRole(...)")}} para permisos por método. \\
\texttt{\seqsplit{io.jsonwebtoken:jjwt-\allowbreak api/impl/jackson}} (0.11.5) & Maven Central, \texttt{\seqsplit{io.jsonwebtoken}} & Emisión y \textbf{verificación de firma} de tokens JWT (HS256) --- validación de autenticidad e integridad del lado del servidor. \\
\texttt{\seqsplit{spring-boot-starter-\allowbreak data-jpa}} (3.4.0) & Maven Central, \texttt{\seqsplit{org.springframework.boot}} & Acceso a datos vía Hibernate ORM con consultas derivadas y parámetros ligados (\emph{prepared statements}), previniendo inyección SQL. \\
\texttt{\seqsplit{com.fasterxml.jackson}} (transitiva, vía \texttt{\seqsplit{spring-boot-starter-web}}) & Maven Central, \texttt{\seqsplit{com.fasterxml.jackson.core}} & Serialización/deserialización JSON con tipado estricto (rechaza tipos incompatibles con \texttt{\seqsplit{InvalidFormatException}}) y escapado automático de caracteres especiales en la salida JSON. \\
\texttt{\seqsplit{org.postgresql:postgresql}} (runtime) & Maven Central, \texttt{\seqsplit{org.postgresql}} & Driver JDBC; ejecuta las sentencias preparadas generadas por Hibernate. \\
\texttt{\seqsplit{org.flywaydb:flyway-core}} & Maven Central, \texttt{\seqsplit{org.flywaydb}} & Migraciones de esquema versionadas (integridad estructural de los datos persistidos). \\
\texttt{\seqsplit{springdoc-openapi-\allowbreak starter-webmvc-ui}} (2.8.5) & Maven Central, \texttt{\seqsplit{org.springdoc}} & Genera la documentación OpenAPI/Swagger a partir de las mismas anotaciones de validación (contrato de entrada siempre sincronizado con el código). \\
\bottomrule
\end{longtable}

\noindent\textbf{Decisión de diseño.} No se usa ninguna librería de sanitización de HTML/JS (p. ej. OWASP Java HTML Sanitizer o ESAPI) porque el backend \textbf{no renderiza HTML}: es una API JSON pura consumida por clientes que hacen su propio \emph{output encoding} al presentar los datos.

\section{Backend: Desarrollo --- apartados validados (con código y evidencia)}

\subsection{2. Validación del lado del servidor --- \Cumple}

\textbf{Estado:} \Cumple. Las cuatro subcategorías del checklist institucional (\emph{autenticidad}, \emph{consistencia}, \emph{integridad}, \emph{permisos}) están cubiertas:

\begin{itemize}[leftmargin=1.2cm]
    \item \textbf{Autenticidad e integridad del token:} \texttt{\seqsplit{JwtUtil}} firma los tokens con HMAC-SHA256 y, al validarlos, \texttt{\seqsplit{Jwts.parserBuilder().setSigningKey(...).parseClaimsJws(token)}} recalcula y verifica la firma; cualquier alteración del payload o de la firma lanza \texttt{\seqsplit{SecurityException}}/\texttt{\seqsplit{MalformedJwtException}}, y la expiración se comprueba aparte.
    \item \textbf{Consistencia:} las referencias a catálogos enviadas por el cliente (perfil, municipio, estado) se verifican contra la base de datos antes de persistir (\texttt{\seqsplit{ensureExists}}).
    \item \textbf{Permisos:} cada endpoint mutante está anotado con \texttt{\seqsplit{@PreAuthorize("hasRole('ADMIN')")}}, evaluado por Spring Security con \texttt{\seqsplit{@EnableMethodSecurity}}.
\end{itemize}

\textbf{Dónde se cumple:}
\begin{itemize}[leftmargin=1.2cm,nosep]
\item \texttt{\seqsplit{JwtUtil.java}}
\item \texttt{\seqsplit{JwtAuthenticationFilter.java}}
\item \texttt{\seqsplit{SecurityConfig.java}}
\item \texttt{\seqsplit{UserServiceImpl.java}}
\end{itemize}

\begin{lstlisting}[style=javastyle,caption={Verificación criptográfica de la firma y expiración del JWT (JwtUtil.java)}]
private Claims getAllClaimsFromToken(String token) {
    return Jwts.parserBuilder()
            .setSigningKey(getSigningKey())
            .build()
            .parseClaimsJws(token)   // valida firma; lanza excepcion si fue alterado
            .getBody();
}

public Boolean validateToken(String token, UserDetails userDetails) {
    try {
        final String username = getUsernameFromToken(token);
        return (username.equals(userDetails.getUsername()) && !isTokenExpired(token));
    } catch (SecurityException | MalformedJwtException e) {
        logger.error("Firma JWT invalida: {}", e.getMessage());
    } catch (ExpiredJwtException e) {
        logger.error("El token JWT ha expirado: {}", e.getMessage());
    }
    return false;
}
\end{lstlisting}

\begin{lstlisting}[style=javastyle,caption={Autorización por rol a nivel de método (UserController.java)}]
@GetMapping
@PreAuthorize("hasRole('ADMIN')")
public ResponseEntity<List<UserResponse>> getAllUsers() { /* ... */ }

@DeleteMapping("/{id}")
@PreAuthorize("hasRole('ADMIN')")
public ResponseEntity<Void> deleteUser(@PathVariable Long id) { /* ... */ }
\end{lstlisting}

\begin{lstlisting}[style=javastyle,caption={Reglas de autorización por ruta (SecurityConfig.java)}]
.authorizeHttpRequests(auth -> auth
        .requestMatchers("/v1/auth/admin", "/v1/auth/student").permitAll()
        .requestMatchers("/v1/users/register").permitAll()
        .requestMatchers("/api/actuator/health", "/v1/health").permitAll()
        .anyRequest().authenticated()
);
http.addFilterBefore(jwtAuthenticationFilter, UsernamePasswordAuthenticationFilter.class);
\end{lstlisting}

\subsection{3. Validación de tipo --- \Cumple}

\textbf{Estado:} \Cumple. Todos los cuerpos de petición se deserializan a \emph{records} de Java con campos fuertemente tipados (\texttt{\seqsplit{String}}, \texttt{\seqsplit{Integer}}, \texttt{\seqsplit{Boolean}}, \texttt{\seqsplit{Long}}); si el JSON entrante trae un tipo incompatible (por ejemplo, una cadena donde se espera un booleano), Jackson lanza \texttt{\seqsplit{InvalidFormatException}}, capturada por \texttt{\seqsplit{GlobalExceptionHandler}} y traducida a un HTTP 400 con el nombre del campo y el tipo esperado.

\textbf{Dónde se cumple:}
\begin{itemize}[leftmargin=1.2cm,nosep]
\item \texttt{\seqsplit{GlobalExceptionHandler.java}}
\end{itemize}

\begin{lstlisting}[style=javastyle,caption={Manejo de tipo de dato incorrecto en el cuerpo JSON (GlobalExceptionHandler.java)}]
@ExceptionHandler(HttpMessageNotReadableException.class)
public ResponseEntity<ErrorResponse> handleHttpMessageNotReadable(HttpMessageNotReadableException ex) {
    String detailMsg = "El formato de los datos enviados es incorrecto o el JSON esta mal formado.";
    Throwable cause = ex.getCause();
    if (cause instanceof InvalidFormatException ife) {
        if (!ife.getPath().isEmpty()) {
            String fieldName = ife.getPath().get(0).getFieldName();
            String targetType = ife.getTargetType().getSimpleName();
            detailMsg = "Tipo de dato incorrecto, el campo '" + fieldName
                    + "' debe ser de tipo " + targetType + ".";
        }
    }
    ErrorItem errorItem = ErrorItem.builder()
            .status("400").code("ERR_BAD_REQUEST_FORMAT")
            .title("Bad Request Format").detail(detailMsg).build();
    return new ResponseEntity<>(ErrorResponse.builder().errors(List.of(errorItem)).build(),
            HttpStatus.BAD_REQUEST);
}
\end{lstlisting}

\subsection{4. Lógica de negocio --- \Cumple}

\textbf{Estado:} \Cumple. \texttt{\seqsplit{UserServiceImpl}} aplica las reglas de negocio del dominio antes de persistir: un estudiante no puede registrarse dos veces con el mismo usuario/correo/CURP, y toda referencia a un catálogo (perfil, municipio, estado) debe existir.

\textbf{Dónde se cumple:}
\begin{itemize}[leftmargin=1.2cm,nosep]
\item \texttt{\seqsplit{UserServiceImpl.java}}
\end{itemize}

\begin{lstlisting}[style=javastyle,caption={Reglas de negocio de registro de estudiante (UserServiceImpl.java)}]
public StudentResponse registerStudent(StudentRequest request) {
    String username = request.person().curp().toLowerCase();

    if (userRepository.existsByUsername(username)) {
        throw new BadRequestException("Student is already registered (username already exists).");
    }
    if (userRepository.existsByEmail(request.email())) {
        throw new BadRequestException("El correo electronico ya esta registrado.");
    }
    if (personRepository.existsByCurp(request.person().curp())) {
        throw new BadRequestException("El CURP proporcionado ya esta registrado en el sistema.");
    }
    if (request.profile() != null) {
        ensureExists(profileRepository, request.profile().longValue(), "Perfil");
    }
    if (request.address() != null && request.address().municipalityId() != null) {
        ensureExists(municipalityRepository, request.address().municipalityId().longValue(), "Municipio");
    }
    // ...
}

private <T> void ensureExists(JpaRepository<T, Long> repository, Long id, String entityName) {
    if (id != null && !repository.existsById(id)) {
        throw new ResourceNotFoundException(
            "El " + entityName + " proporcionado (" + id + ") no existe en el catalogo.");
    }
}
\end{lstlisting}

\subsection{8c. Validación de tipo de datos (estructuras JSON) --- \Cumple}

\textbf{Estado:} \Cumple. La validación en cascada de \texttt{\seqsplit{@Valid}} permite que estructuras JSON anidadas (objeto \texttt{\seqsplit{person}} dentro de \texttt{\seqsplit{StudentRequest}}, con su propia lista de \texttt{\seqsplit{relatives}}) se validen recursivamente en un único paso declarativo, sin código manual de recorrido.

\textbf{Dónde se cumple:}
\begin{itemize}[leftmargin=1.2cm,nosep]
\item \texttt{\seqsplit{StudentRequest.java}}
\end{itemize}

\begin{lstlisting}[style=javastyle,caption={Validación en cascada sobre estructuras JSON anidadas (StudentRequest.java)}]
public record StudentRequest(

    @Valid
    @NotNull(message = "Personal data is mandatory")
    PersonRequest person,

    @Valid
    @NotNull(message = "Address is mandatory")
    AddressRequest address,

    @Valid
    @NotNull(message = "Relatives list is mandatory")
    List<RelativeRequest> relatives,

    @NotNull(message = "Profile is mandatory")
    Integer profile,

    @NotBlank(message = "Email is mandatory")
    String email
) {}
\end{lstlisting}

\subsection{8f. Uso de funciones y librerías seguras (ORM) --- \Cumple}

\textbf{Estado:} \Cumple. Ningún repositorio del proyecto usa \texttt{\seqsplit{@Query}} con SQL nativo ni concatenación de cadenas: todas las consultas son métodos derivados de \texttt{\seqsplit{JpaRepository}}, resueltos por Hibernate en sentencias preparadas con parámetros ligados. Esto elimina la clase de inyección SQL clásica en la capa de acceso a datos tal como está escrita hoy.

\textbf{Dónde se cumple:}
\begin{itemize}[leftmargin=1.2cm,nosep]
\item \texttt{\seqsplit{PersonRepository.java}}
\item \texttt{\seqsplit{UserRepository.java}}
\end{itemize}

\begin{lstlisting}[style=javastyle,caption={Consultas 100 porciento parametrizadas vía Spring Data JPA (PersonRepository.java)}]
@Repository
public interface PersonRepository extends JpaRepository<Person, Long> {
    boolean existsByCurp(String curp);   // -> SELECT ... WHERE curp = ?  (parametro ligado)
    Optional<Person> findByCurp(String curp);
}
\end{lstlisting}

\subsection{8i. Escape de salida contextual (JSON) --- \Cumple}

\textbf{Estado:} \Cumple. Toda respuesta se serializa con Jackson (\texttt{\seqsplit{com.fasterxml.jackson.databind.ObjectMapper}}), que escapa automáticamente comillas, barras invertidas y caracteres de control al producir el cuerpo JSON de salida, evitando que un valor persistido con comillas o caracteres especiales rompa la estructura de la respuesta. No hay generación manual de HTML/JS en el servidor, por lo que no aplica el escapado de esos contextos (ver Sección~\ref{sec:backend-noaplica}).

\textbf{Dónde se cumple:} todos los \texttt{\seqsplit{@RestController}} del proyecto (serialización automática vía \texttt{\seqsplit{spring-boot-starter-web}} / Jackson).

\subsection{9. Uso de librerías y frameworks de validación --- \Cumple}

\textbf{Estado:} \Cumple. Toda la validación declarativa usa \textbf{Jakarta Bean Validation}, un estándar de la plataforma Java (JSR 380) implementado por \textbf{Hibernate Validator}, ambos mantenidos activamente y ampliamente auditados; no existe ningún validador propio ni ad-hoc en el backend (a diferencia del cliente, que sí construyó un módulo propio por no depender de librerías externas en la lógica de seguridad de una app sin servidor).

\textbf{Dónde se cumple:}
\begin{itemize}[leftmargin=1.2cm,nosep]
\item \texttt{\seqsplit{pom.xml}}
\item DTOs listados en la Sección~\ref{ssec:backend-tipos}
\end{itemize}

\begin{lstlisting}[style=xmlstyle,caption={Dependencia que habilita Jakarta Bean Validation / Hibernate Validator (pom.xml)}]
<dependency>
    <groupId>org.springframework.boot</groupId>
    <artifactId>spring-boot-starter-validation</artifactId>
</dependency>
\end{lstlisting}

\subsection{11. Gestión de errores adecuada --- \Cumple}

\textbf{Estado:} \Cumple. \texttt{\seqsplit{GlobalExceptionHandler}} centraliza el manejo de excepciones y evita que el servidor filtre información sensible: los errores de validación (\texttt{\seqsplit{MethodArgumentNotValidException}}) devuelven \textbf{422} con un mensaje por campo; las violaciones de integridad de datos (\texttt{\seqsplit{DataIntegrityViolationException}}) devuelven \textbf{409} con un mensaje genérico (sin exponer la sentencia SQL ni el nombre de la restricción); y cualquier excepción no controlada cae en un manejador \texttt{\seqsplit{Exception.class}} que registra el detalle solo en el log del servidor y responde \textbf{500} con un mensaje genérico.

\textbf{Dónde se cumple:}
\begin{itemize}[leftmargin=1.2cm,nosep]
\item \texttt{\seqsplit{GlobalExceptionHandler.java}}
\end{itemize}

\begin{lstlisting}[style=javastyle,caption={Traducción de errores de validación a HTTP 422, uno por campo (GlobalExceptionHandler.java)}]
@ExceptionHandler(MethodArgumentNotValidException.class)
public ResponseEntity<ErrorResponse> handleValidationErrors(MethodArgumentNotValidException ex) {
    List<ErrorItem> errorItems = ex.getBindingResult().getFieldErrors().stream()
            .map(error -> ErrorItem.builder()
                    .status("422")
                    .code("ERR_VAL_" + error.getField().toUpperCase())
                    .title("Unprocessable Entity")
                    .detail(error.getDefaultMessage())
                    .source(ErrorSource.builder()
                            .pointer("/data/attributes/" + error.getField()).build())
                    .build())
            .collect(Collectors.toList());
    return new ResponseEntity<>(ErrorResponse.builder().errors(errorItems).build(),
            HttpStatus.UNPROCESSABLE_ENTITY);
}

@ExceptionHandler(Exception.class)
public ResponseEntity<ErrorResponse> handleGeneralException(Exception ex) {
    logger.error("Error no controlado del sistema: ", ex);  // detalle solo en el log del servidor
    ErrorItem errorItem = ErrorItem.builder()
            .status("500").code("ERR_INTERNAL_SERVER").title("Internal Server Error")
            .detail("Ocurrio un error interno en el servidor. Por favor, intente mas tarde.")
            .build();
    return new ResponseEntity<>(ErrorResponse.builder().errors(List.of(errorItem)).build(),
            HttpStatus.INTERNAL_SERVER_ERROR);
}
\end{lstlisting}

\section{Backend: Apartados Parciales}
\label{sec:backend-parcial}

\subsection{5. Patrones y reglas específicas --- \Parcial}

\textbf{Estado:} \Parcial. El correo electrónico se valida con \texttt{\seqsplit{@Email}} en \texttt{\seqsplit{RegisterRequest}} (\texttt{\seqsplit{@NotBlank}}, \texttt{\seqsplit{@Email}}, \texttt{\seqsplit{@Size}}), y la contraseña valida longitud mínima y máxima.

\textbf{Dónde se cumple:}
\begin{itemize}[leftmargin=1.2cm,nosep]
\item \texttt{\seqsplit{RegisterRequest.java}}
\end{itemize}

\textbf{Por qué es parcial:} \texttt{\seqsplit{StudentRequest.email}} solo tiene \texttt{\seqsplit{@NotBlank}} (sin \texttt{\seqsplit{@Email}}), por lo que un correo mal formado puede pasar por la ruta de registro de estudiante; y la contraseña no exige \texttt{\seqsplit{@Pattern}} de mezcla de mayúsculas/números/símbolos, solo longitud (\texttt{\seqsplit{@Size(min = 6, max = 100)}}). No aplica Luhn de tarjetas porque el dominio no maneja pagos.

\textbf{Recomendación:} añadir \texttt{\seqsplit{@Email}} a \texttt{\seqsplit{StudentRequest.email}} y una \texttt{\seqsplit{@Pattern}} de complejidad mínima de contraseña.

\subsection{7. Validación contextual --- \Parcial}

\textbf{Estado:} \Parcial. \texttt{\seqsplit{ensureExists}} valida que un \texttt{\seqsplit{municipalityId}} y un \texttt{\seqsplit{stateId}} existan en sus catálogos antes de persistir.

\textbf{Dónde se cumple:}
\begin{itemize}[leftmargin=1.2cm,nosep]
\item \texttt{\seqsplit{UserServiceImpl.java}}
\end{itemize}

\textbf{Por qué es parcial:} la verificación es \emph{individual} por campo; no comprueba que el municipio efectivamente pertenezca al estado enviado (validación cruzada contextual entre ambos campos).

\textbf{Recomendación:} extender \texttt{\seqsplit{ensureExists}} o agregar una comprobación específica \texttt{\seqsplit{municipalityBelongsToState(municipalityId, stateId)}}.

\subsection{8b. Filtrado (allowlist/blacklist) --- \Parcial}

\textbf{Estado:} \Parcial. Los DTO usan Bean Validation para exigir formato en algunos campos (\texttt{\seqsplit{@Email}}).

\textbf{Dónde se cumple:}
\begin{itemize}[leftmargin=1.2cm,nosep]
\item \texttt{\seqsplit{RegisterRequest.java}}
\end{itemize}

\textbf{Por qué es parcial:} no hay una \texttt{\seqsplit{@Pattern}} de \emph{allowlist} de caracteres para campos de texto libre como \texttt{\seqsplit{username}}, \texttt{\seqsplit{nombre}} de catálogos o \texttt{\seqsplit{descripcion}} de rol (columna \texttt{\seqsplit{TEXT}} sin restricción), a diferencia del cliente que sí define un allowlist ASCII explícito (\texttt{\seqsplit{printableAscii}}).

\textbf{Recomendación:} agregar \texttt{\seqsplit{@Pattern}} a \texttt{\seqsplit{username}} y evaluar longitud máxima para \texttt{\seqsplit{RoleRequestDTO.descripcion}}.

\subsection{8d. Limpieza de entradas (trim/normalize) --- \Parcial}

\textbf{Estado:} \Parcial. Existe normalización puntual: \texttt{\seqsplit{request.person().curp().toLowerCase()}} construye el \texttt{\seqsplit{username}} en minúsculas.

\textbf{Dónde se cumple:}
\begin{itemize}[leftmargin=1.2cm,nosep]
\item \texttt{\seqsplit{UserServiceImpl.java}}
\end{itemize}

\textbf{Por qué es parcial:} \texttt{\seqsplit{@NotBlank}} rechaza cadenas vacías o compuestas solo por espacios, pero \textbf{no recorta} espacios al inicio/fin de una cadena válida antes de persistirla (por ejemplo, \texttt{\seqsplit{" admin "}} pasa la validación y se guarda con espacios).

\textbf{Recomendación:} aplicar \texttt{\seqsplit{.strip()}} a los campos de texto en el servicio, o un deserializador Jackson que recorte espacios globalmente.

\subsection{8j. Revisiones y auditorías (SAST/pentest) --- \Parcial}

\textbf{Estado:} \Parcial. El código se revisó manualmente para este informe (anotaciones de validación, manejo de excepciones, configuración de seguridad).

\textbf{Por qué es parcial:} a diferencia del cliente (que sí tiene \texttt{\seqsplit{.github/workflows/security.yml}} con \texttt{\seqsplit{flutter analyze}}, \texttt{\seqsplit{osv-scanner}} y \texttt{\seqsplit{gitleaks}}), el backend \textbf{no tiene todavía} pipeline de CI ni herramienta de análisis estático (no hay \texttt{\seqsplit{spotbugs}}, \texttt{\seqsplit{checkstyle}} ni \texttt{\seqsplit{sonar}} configurados en \texttt{\seqsplit{pom.xml}}), y \texttt{\seqsplit{src/test/java}} está vacío --- no hay pruebas unitarias ni de integración que verifiquen automáticamente estas reglas de validación.

\textbf{Recomendación:} agregar un workflow de GitHub Actions equivalente al del cliente (\texttt{\seqsplit{mvn verify}}, \texttt{\seqsplit{spotbugs}}, \texttt{\seqsplit{osv-scanner}}, \texttt{\seqsplit{gitleaks}}) y una primera batería de pruebas \texttt{\seqsplit{@WebMvcTest}}/\texttt{\seqsplit{MockMvc}} que ejerciten los DTO documentados en la Sección~\ref{ssec:backend-tipos}.

\section{Backend: Apartados No aplica (justificados)}
\label{sec:backend-noaplica}

\begin{description}[leftmargin=1.4cm,style=nextline]
\item[1. Validación del lado del cliente] \textbf{Estado:} \NoAplica\ para este componente. \textbf{Por qué no aplica:} es responsabilidad exclusiva del cliente móvil, ya documentada en la Sección~\ref{sec:cliente-desarrollo}.
\item[6. Validación cruzada (fechas/campos generales)] \textbf{Estado:} \NoAplica. \textbf{Por qué no aplica:} ningún DTO actual expone pares de campos tipo "fecha inicio/fecha fin"; el único caso de coherencia entre campos identificado (municipio/estado) se documenta como \Parcial\ en la Sección~\ref{sec:backend-parcial}.
\item[8a. Escapado HTML/JS/SQL] \textbf{Estado:} \NoAplica. \textbf{Por qué no aplica:} el backend no genera HTML ni JavaScript (es una API JSON pura); el escapado SQL se resuelve estructuralmente vía JPA/Hibernate (ver 8f, \Cumple), no por escapado manual de caracteres.
\item[8e. Codificación de entradas (Base64/URL)] \textbf{Estado:} \NoAplica. \textbf{Por qué no aplica:} ningún endpoint recibe payloads binarios en Base64; la decodificación URL de parámetros de ruta/consulta la resuelve el propio contenedor de servlets de forma transparente.
\item[8g. Reemplazo de caracteres] \textbf{Estado:} \NoAplica. \textbf{Por qué no aplica:} no hay generación de HTML/SQL manual que requiera sustituir comillas o etiquetas de script.
\item[8h. Canonicalización de rutas de archivo] \textbf{Estado:} \NoAplica. \textbf{Por qué no aplica:} el backend no construye rutas de sistema de archivos a partir de datos de usuario; el único acceso a \texttt{\seqsplit{classpath}} (\texttt{\seqsplit{FirebaseConfig}}) usa un recurso fijo y no controlado por el cliente.
\item[10. Educación y capacitación] \textbf{Estado:} \NoAplica. \textbf{Por qué no aplica:} no evidenciable desde el repositorio (igual criterio que en el documento del cliente).
\end{description}

\section{Backend: Hallazgos de seguridad y recomendaciones}

Durante la elaboración de este informe se identificaron los siguientes hallazgos, adicionales a los ya cubiertos como apartados \Parcial. Se documentan con severidad orientativa y no implican necesariamente una corrección ya aplicada.

\begin{longtable}{p{4.9cm}p{2cm}p{6.3cm}}
\toprule
\textbf{Hallazgo} & \textbf{Severidad} & \textbf{Detalle / recomendación} \\
\midrule
\endhead
Campos sin cota de longitud alineada con la base de datos (\texttt{\seqsplit{curp}}, \texttt{\seqsplit{zipCode}}, \texttt{\seqsplit{imageUrl}}) & Media & \texttt{\seqsplit{PersonRequest.curp}} no tiene \texttt{\seqsplit{@Size(max=18)}} pese a que la columna es \texttt{\seqsplit{VARCHAR(18)}}; una entrada más larga falla como \texttt{\seqsplit{DataIntegrityViolationException}} (409 genérico) en vez de un 422 claro por campo. Agregar \texttt{\seqsplit{@Size}} espejo de cada \texttt{\seqsplit{@Column(length=...)}}. \\
\texttt{\seqsplit{ApiKeyFilter}} inyectado pero no registrado en la cadena de seguridad & Media & \texttt{\seqsplit{SecurityConfig}} declara \texttt{\seqsplit{@Autowired ApiKeyFilter}} pero nunca lo añade con \texttt{\seqsplit{http.addFilterBefore(...)}}; el filtro igual se activa globalmente por ser un \texttt{\seqsplit{@Component}}/\texttt{\seqsplit{OncePerRequestFilter}} auto-registrado por Spring Boot. Documentar explícitamente esta doble vía de activación o integrarlo formalmente a \texttt{\seqsplit{SecurityFilterChain}}. \\
Clave de API no está realmente hasheada & Media & El campo \texttt{\seqsplit{keyHash}} y el método \texttt{\seqsplit{findByKeyHash}} sugieren un hash, pero se compara y almacena el valor generado (\texttt{\seqsplit{UUID.randomUUID()}} concatenado) en texto plano. Aplicar una función de hash (p. ej. SHA-256) antes de persistir y comparar. \\
Valores por defecto de secretos en \texttt{\seqsplit{application.yml}} & Media & \texttt{\seqsplit{JWT\_SECRET}}, la contraseña de base de datos y \texttt{\seqsplit{ENCRYPTION\_SECRET}} tienen valores de respaldo con apariencia real hardcodeados en el archivo versionado. Aunque están pensados para sobrescribirse por variable de entorno, se recomienda usar un gestor de secretos y no comprometer valores plausibles en el repositorio. \\
Sin límite de tasa (\emph{rate limiting}) ni CAPTCHA en endpoints públicos & Baja/Media & \texttt{\seqsplit{/v1/auth/admin}}, \texttt{\seqsplit{/v1/auth/student}} y \texttt{\seqsplit{/v1/users/register}} son públicos y no tienen control de fuerza bruta ni de registro masivo. Evaluar Bucket4j/resilience4j o un límite a nivel de proxy. \\
Sin cota de tamaño de \texttt{\seqsplit{Pageable.size}} en peticiones GET & Baja & No se configuró \texttt{\seqsplit{spring.data.web.pageable.max-page-size}}; aplica el máximo por defecto de Spring Boot (2000), lo que permite solicitar páginas grandes. Fijar un máximo explícito acorde al catálogo más grande. \\
\bottomrule
\end{longtable}

\section{Backend: Tabla Resumen de Cumplimiento}

\begin{longtable}{p{0.5cm}p{4cm}p{2.6cm}p{6.2cm}}
\toprule
\textbf{\#} & \textbf{Apartado} & \textbf{Estado} & \textbf{Evidencia} \\
\midrule
\endhead
1 & Validación cliente & \NoAplica & Responsabilidad del cliente móvil \\
2 & Validación servidor & \Cumple & \texttt{\seqsplit{JwtUtil.java}}, \texttt{\seqsplit{SecurityConfig.java}} \\
3 & Validación de tipo & \Cumple & \texttt{\seqsplit{GlobalExceptionHandler.java}} (Jackson) \\
4 & Lógica de negocio & \Cumple & \texttt{\seqsplit{UserServiceImpl.java}} \\
5 & Patrones (email/contraseña) & \Parcial & \texttt{\seqsplit{RegisterRequest.java}} vs. \texttt{\seqsplit{StudentRequest.java}} \\
6 & Validación cruzada & \NoAplica & Sin pares fecha inicio/fin en el dominio \\
7 & Validación contextual & \Parcial & \texttt{\seqsplit{ensureExists}} (sin cruce municipio/estado) \\
8a & Escapado HTML/JS/SQL & \NoAplica & API JSON pura; SQL resuelto por JPA \\
8b & Filtrado (allowlist) & \Parcial & Sin \texttt{\seqsplit{@Pattern}} en texto libre \\
8c & Validación tipo JSON & \Cumple & \texttt{\seqsplit{StudentRequest.java}} (cascada \texttt{\seqsplit{@Valid}}) \\
8d & Limpieza (trim/normalize) & \Parcial & \texttt{\seqsplit{@NotBlank}} no recorta espacios \\
8e & Codificación & \NoAplica & Sin payloads Base64; URL decoding del contenedor \\
8f & Librerías seguras / ORM & \Cumple & \texttt{\seqsplit{PersonRepository.java}} (JPA parametrizado) \\
8g & Reemplazo de caracteres & \NoAplica & Sin generación manual de HTML/SQL \\
8h & Canonicalización & \NoAplica & Sin manejo de rutas de archivo de usuario \\
8i & Escape de salida & \Cumple & Serialización JSON vía Jackson \\
8j & Revisiones/auditorías & \Parcial & Sin CI ni pruebas automatizadas todavía \\
9 & Librerías de validación & \Cumple & \texttt{\seqsplit{spring-boot-starter-validation}} \\
10 & Educación/capacitación & \NoAplica & No evidenciable desde el repositorio \\
11 & Gestión de errores & \Cumple & \texttt{\seqsplit{GlobalExceptionHandler.java}} \\
\bottomrule
\end{longtable}

\noindent\textbf{Conteo final (20 apartados, Backend):} \Cumple\ = \textbf{8} (2, 3, 4, 8c, 8f, 8i, 9, 11) \quad \Parcial\ = \textbf{5} (5, 7, 8b, 8d, 8j) \quad \NoAplica\ = \textbf{7} (1, 6, 8a, 8e, 8g, 8h, 10) \quad \Fuera\ = 0 \quad \NoCumple\ = 0.

\section{Tabla Consolidada: Cliente vs. Backend}

La siguiente tabla cruza el estado de cada apartado en ambos componentes, mostrando cómo el backend cierra la mayoría de los apartados que el cliente había dejado \Fuera\ de su alcance.

\begin{longtable}{p{0.5cm}p{4cm}p{2.3cm}p{2.3cm}p{5cm}}
\toprule
\textbf{\#} & \textbf{Apartado} & \textbf{Cliente} & \textbf{Backend} & \textbf{Lectura combinada} \\
\midrule
\endhead
1 & Validación cliente & \Cumple & \NoAplica & Cubierto en el cliente \\
2 & Validación servidor & \Fuera & \Cumple & Cerrado por el backend (JWT + roles) \\
3 & Validación de tipo & \Cumple & \Cumple & Cubierto en ambas capas \\
4 & Lógica de negocio & \Cumple & \Cumple & Cubierto en ambas capas \\
5 & Patrones específicos & \NoAplica & \Parcial & Solo relevante en backend; falta uniformar \texttt{\seqsplit{@Email}} \\
6 & Validación cruzada & \NoAplica & \NoAplica & No aplica en ninguna capa hoy \\
7 & Validación contextual & \Cumple & \Parcial & Backend pendiente de cruce municipio/estado \\
8a & Escapado HTML/JS/SQL & \NoAplica & \NoAplica & No aplica en ninguna capa \\
8b & Filtrado (allowlist) & \Cumple & \Parcial & Backend sin \texttt{\seqsplit{@Pattern}} de texto libre \\
8c & Validación tipo JSON & \Cumple & \Cumple & Cubierto en ambas capas \\
8d & Limpieza (trim/normalize) & \Cumple & \Parcial & Backend no recorta espacios \\
8e & Codificación & \NoAplica & \NoAplica & No aplica en ninguna capa \\
8f & Librerías seguras / ORM & \Cumple & \Cumple & Cubierto en ambas capas \\
8g & Reemplazo de caracteres & \NoAplica & \NoAplica & No aplica en ninguna capa \\
8h & Canonicalización & \Cumple & \NoAplica & Solo relevante en cliente (URLs de medios) \\
8i & Escape de salida & \NoAplica & \Cumple & Solo relevante en backend (serialización JSON) \\
8j & Revisiones/auditorías & \Parcial & \Parcial & Pendiente en ambas capas (backend más atrás: sin CI) \\
9 & Librerías de validación & \Cumple & \Cumple & Cubierto en ambas capas \\
10 & Educación/capacitación & \NoAplica & \NoAplica & No evidenciable en ninguna capa \\
11 & Gestión de errores & \Cumple & \Cumple & Cubierto en ambas capas \\
\bottomrule
\end{longtable}

\section{Conclusiones Generales (Cliente + Backend)}

El cliente móvil llegó a un estado conforme: 11 apartados en \Cumple, 1 en \Parcial\ (8j) y 0 en \NoCumple, tras cerrarse las brechas de mayor severidad bajo su control (token cifrado, logs sin secretos, entrada externa validada).

El backend, evaluado por primera vez en este informe, cubre \textbf{8 apartados en \Cumple}, incluyendo exactamente el que el cliente había dejado \Fuera\ de su alcance (validación del lado del servidor, mediante verificación criptográfica de JWT y autorización por rol). Quedan \textbf{5 apartados en \Parcial}, todos con una acción concreta y acotada: uniformar \texttt{\seqsplit{@Email}} entre \texttt{\seqsplit{RegisterRequest}} y \texttt{\seqsplit{StudentRequest}}, cruzar municipio/estado, añadir \texttt{\seqsplit{@Pattern}} de allowlist a texto libre, recortar espacios en cadenas, y --- el más relevante para la madurez del proyecto --- incorporar un pipeline de CI con pruebas automatizadas y análisis estático, tal como ya existe en el repositorio del cliente.

\textbf{Riesgos abiertos, externos a la validación de entradas:} (1) TLS/HTTPS en los servicios de evaluación (infraestructura, severidad Alta, ya documentado en la Sección~4.2); (2) la clave de API no está realmente hasheada y se transmite en texto plano vía FCM; (3) valores de secretos con apariencia real en \texttt{\seqsplit{application.yml}}; (4) ausencia de límite de tasa en los endpoints públicos de autenticación y autorregistro.

\textbf{Recomendaciones priorizadas:} (1) agregar pruebas \texttt{\seqsplit{MockMvc}} y un workflow de CI al backend; (2) uniformar las anotaciones de Bean Validation faltantes (\texttt{\seqsplit{@Email}}, \texttt{\seqsplit{@Size}} espejo de \texttt{\seqsplit{@Column}}, \texttt{\seqsplit{@Pattern}} de allowlist); (3) hashear la clave de API y formalizar su registro en \texttt{\seqsplit{SecurityFilterChain}}; (4) mover los secretos de \texttt{\seqsplit{application.yml}} a un gestor dedicado; (5) una vez exista TLS en los servicios de evaluación, aplicar \emph{certificate pinning} en el cliente y forzar HTTPS.

\end{document}
